\subsection{Interaktionsprinzipien nach ISO 9241-110}   % zur Gebrauchstauglichkeit und Ergonomie interaktiver Systeme
\label{ch:SdF:UI:sec:ISO-Interaktionsprinzipien}

Der Text der ISO 9241 Part 110 kombiniert die Ausprägungen der Usability mit Gestaltungsempfehlungen in Form von sieben Interaktionsprinzipien (ehem. Grundsätze der Dialoggestaltung). Indem Systeme und Schnittstellen gebrauchstauglicher, zugänglicher und konsistenter gestaltet werden, erleichtern und verbessern diese die kognitive Ergonomie der Mensch-System-Interaktion.

Die sieben Interaktionsprinzipien sind die Aufgabenangemessenheit, Selbstbeschreibungsfähigkeit, Erwartungskonformität, Erlernbarkeit, Steuerbarkeit, Robustheit gegen Benutzungsfehler und die Benutzerbindung. Die Gewichtung der Grundsätze kann unterschiedlich ausfallen, da diese von dem System, den Nutzern und deren Aufgaben abhängig sind. Zur Bewertung der Umsetzung der Empfehlungen ist die \glqq Tabelle zur Unterstützung der Anwendung [der ISO 9241-110]\grqq{} \cite{DIN_EN_ISO_9241-110} der ISO-Norm beigelegt. Die vollständige Version ist in \autoref{Add:ISO-Tabellen} zu finden.
Anhand der nachfolgenden Beschreibungen werden die Interaktionsprinzipien wie in \autoref{tab:usability<->ISO-Prinzipien} zu sehen, den drei Usability-Dimensionen zugeordnet.

% Please add the following required packages to your document preamble:
% \usepackage[table,xcdraw]{xcolor}
% Beamer presentation requires \usepackage{colortbl} instead of \usepackage[table,xcdraw]{xcolor}
\begin{table}[H]
    % Farben
    % grün 9ADD99 96C800
    % gelb FFFC8B
    \centering
    \begin{small}
    \begin{tabular}{ll|ccc|}
    \cline{3-5}
                                                                                           &                                                                             & \multicolumn{3}{c|}{Usability}                                                                                                                \\ \cline{3-5} 
                                                                                          &                                                                             & \multicolumn{1}{c|}{Effektivität}                    & \multicolumn{1}{c|}{Effizienz}                       & Zufriedenheit                   \\ \hline
    \multicolumn{1}{|c|}{\multirow{7}{*}{\rotatebox{90}{\makecell[c]{Interaktionsprinzipien der \\DIN EN ISO 9241-110}}}} & \begin{tabular}[c]{@{}l@{}}Aufgaben-\\ angemessenheit\end{tabular}          & \multicolumn{1}{c|}{\cellcolor[HTML]{96C800}ja}      & \multicolumn{1}{c|}{\cellcolor[HTML]{96C800}ja}      & nein                            \\ \cline{2-5} 
    \multicolumn{1}{|l|}{}                                                                 & Steuerbarkeit                                                               & \multicolumn{1}{c|}{\cellcolor[HTML]{96C800}ja}      & \multicolumn{1}{c|}{\cellcolor[HTML]{FFFC8B}eher ja} & \cellcolor[HTML]{96C800}ja      \\ \cline{2-5} 
    \multicolumn{1}{|l|}{}                                                                 & \begin{tabular}[c]{@{}l@{}}Erwartungs-\\ konformität\end{tabular}           & \multicolumn{1}{c|}{\cellcolor[HTML]{FFFC8B}eher ja} & \multicolumn{1}{c|}{nein}                            & \cellcolor[HTML]{96C800}ja      \\ \cline{2-5} 
    \multicolumn{1}{|l|}{}                                                                 & Erlernbarkeit                                                               & \multicolumn{1}{c|}{\cellcolor[HTML]{FFFC8B}eher ja} & \multicolumn{1}{c|}{\cellcolor[HTML]{96C800}ja}      & \cellcolor[HTML]{FFFC8B}eher ja \\ \cline{2-5} 
    \multicolumn{1}{|l|}{}                                                                 & \begin{tabular}[c]{@{}l@{}}Selbstbeschreibungs-\\ fähigkeit\end{tabular}    & \multicolumn{1}{c|}{\cellcolor[HTML]{FFFC8B}eher ja} & \multicolumn{1}{c|}{\cellcolor[HTML]{96C800}ja}      & nein                            \\ \cline{2-5} 
    \multicolumn{1}{|l|}{}                                                                 & \begin{tabular}[c]{@{}l@{}}Robustheit gegen\\ Benutzungsfehler\end{tabular} & \multicolumn{1}{c|}{\cellcolor[HTML]{FFFC8B}eher ja} & \multicolumn{1}{c|}{\cellcolor[HTML]{FFFC8B}eher ja} & nein                            \\ \cline{2-5} 
    \multicolumn{1}{|l|}{} & Benutzerbindung                                                             & \multicolumn{1}{c|}{nein}                            & \multicolumn{1}{c|}{nein}                            & \cellcolor[HTML]{96C800}ja      \\ \hline
    \end{tabular}
    \caption{Zuordnung der Interaktionsprinzipien der DIN EN ISO 9241-110 zur den drei Dimensionen der Usability \textcolor{red}{ZUTUN: Farben}}
    \label{tab:usability<->ISO-Prinzipien}
    \end{small}
\end{table}

\noindent
Nachfolgend werden die Inhalte der DIN EN ISO 9241-110 Interaktionsprinzipien \cite{DIN_EN_ISO_9241-110} mit Fokus zur Anwendung auf Dash\-boards zusammengefasst.

% Aufgabenangemessenheit
Ist ein System \textbf{aufgabenangemessen}, bedeutet dies, dass es den Nutzer bei seinen Aufgaben unterstützt. Hierbei ist das System primär auf diese Aufgabe ausgerichtet, anstatt auf der dafür benötigten Technik dahinter zu basieren. Um den Aufwand einer Aufgabe zu optimieren, sollen beispielweise unnötige Schritte sowie ablenkende Informationen oder Interaktionsmöglichkeiten auf der Oberfläche gemieden werden. Dies äußert sich, indem zum Beispiel nur die Informationen auf der Oberfläche angezeigt werden, die für den aktuellen Schritt benötigt werden.

% Selbstbeschreibungsfähigkeit
Die \textbf{Selbstbeschreibungsfähigkeit} hilft dem Nutzer, ein System intuitiv bedienen zu können. Es soll angemessene Informationen und offensichtliche Fähigkeiten des Systems mitteilen, ohne unnötige Interaktionen zu erfordern. Beispiele hierfür sind die Platzierung von Steuerelementen, deren eindeutige Unterscheidung von nicht interaktiven Schnittstellen und dem offensichtlichen Aufzeigen des aktuellen Navigationsverlaufs in Form von etwa breadcrumb trails (dt. Brotkrumenpfade). Diese Maßnahmen können die Notwendigkeit externer Hilfe, wie beispielsweise aus Handbüchern oder von Beratungsstellen, stark reduzieren.

% Erwartungskonformität
Die \textbf{Erwartungskonformität} beschreibt, dass das Verhalten des interaktiven Systems vorhersehbar und somit basierend auf dem aktuellen Kontext mit den Nutzererwartungen konform sein soll.
Es wird erwartet, dass sich das System angemessen verhält und entsprechend reagiert. Dies wird durch Arbeitsschritte gestützt, die mit dem Aufgabenverständnis des Nutzers sowie den Erfordernissen des breitesten Benutzerspektrums übereinstimmen. Dafür ist eine dem Nutzer vertraute Darstellung, Eingabe und Steuerung erforderlich, welche innerhalb des Systems und darüber hinaus konsistent ist.
Das System sollte mit Änderungen im Nutzungskontext zurechtkommen. Dafür soll es in der Lage sein, auf unterschiedliche Erfordernisse einzelner Nutzer sowie auf angeschlossene Ressourcen angemessen zu reagieren. Informationen sind auf verschiedenen Gerätearten und Anzeigegrößen angemessen darzustellen und die Änderungen hinsichtlich der physischen Umgebung des Benutzers anzupassen, wie beispielsweise das Übernehmen der Light- und Darkmode"=Einstellung.

% Erlernbarkeit
Damit ein System \textbf{erlernbar} ist, hat es das Explorieren und die Entdeckung seiner Fähigkeiten und deren Verwendung zu ermöglichen, um die vom Nutzer angestrebten Ergebnisse zu erreichen und ihm den Zugang zu weiteren Informationen zu ermöglichen. Es hat den Lernaufwand zu minimieren und Unterstützung beim Erlernen anzubieten. Dafür sollte es durch das Ausprobieren keine negativen Konsequenzen für den Nutzer geben. Zudem wird die Implementierung geeigneter Alternativen zur Suche nach und für das Navigieren zu Informationen und Funktionen empfohlen.
Weiterhin soll das System Rückmeldung geben, um dem Nutzer zu helfen, die Folgen seiner Aktionen zu erlernen, und in dem Zuge die Entwicklung seiner Fertigkeiten im Zusammenhang mit den unterstützten Aufgaben fördern.

% Steuerbarkeit
Die \textbf{Steuerbarkeit} erlaubt es dem Benutzer, die Kontrolle über die Geschwindigkeit, Abfolge und Individualisierung der Nutzung eines Systems zu behalten. Das beinhaltet die Möglichkeit zur Flexibilität des Nutzers, Aufgaben in einer beliebigen Reihenfolge durchzuführen, diese zu jeder Zeit zu unterbrechen und zu einem späteren Zeitpunkt an derselben Stelle wieder aufnehmen zu können.

% Robustehit gegen Benuzungsfehler
Das System hat eine \textbf{Robustheit gegen Benutzungsfehler} aufzuweisen, indem diese dem Nutzer bei der Fehlervermeidung hilft, eine Toleranz und Hilfestellung im Fehlerfall aufweist sowie den Nutzer dabei unterstützt, diesen Fehler zu beheben.

% Benutzerbindung (en: user engagement)
Das letzte der Prinzipien ist die \textbf{Benutzerbindung}. Diese besagt, dass das System Funktionen und Informationen auf einladende und motivierende Weise darzustellen hat, um eine kontinuierliche Interaktion mit dem System zu fördern und eine positive User-Experience herbeizuführen. Die Benutzerbindung ist aus dem englischen \textit{user engagement} übersetzt und kann ebenfalls als \textit{Benutzer-Engagement} verstanden werden. In der deutschen Sprache ist das Wort \textit{Engagement} im  Vergleich zum Wort \textit{Bindung} oft positiv besetzt, daher kann der englische Begriff \textit{engagement} durch eine Kombination der beiden Bedeutungen beschrieben werden, ist jedoch noch umfassender zu verstehen.
Das System soll den Nutzer motivieren, mit diesem zu interagieren. Dabei werden seine unterbewussten Bedürfnisse antizipiert. Es soll ein Vertrauensverhältnis aufgebaut werden, indem ein konsistentes Design und ein hochwertiges Erscheinungsbild gewahrt wird. Weiterhin soll der Nutzer zunehmend einbezogen werden, indem ihm die Gelegenheit geboten wird, Vorschläge zu Änderungen, Systemerweiterungen und -verbesserungen einzureichen. 

% Allgemein
Die Reihenfolge der Prinzipien stellt keine Rangfolge dar. Die Prinzipien und Empfehlungen sind generischer Natur und damit nicht an ein bestimmtes System oder Kontext gebunden. Der Inhalt der Prinzipien überschneidet sich teilweise, wobei die Bedeutung und Möglichkeit zur Anwendung des jeweiligen Prinzips vom Nutzungskontext, den Benutzergruppen und der gewählten Technik abhängt. 
[Zusammenfassung nach \cite{DIN_EN_ISO_9241-110}]